\section{Dynamics: candidate routes to stabilization}
\label{sec:dynamics}

The dynamics module explains trajectories of construction-level licensing
posteriors, and in principle refinements of the conditioning partition itself.
It doesn't define grammaticality.

\subsection{Niches, competitors, and opportunity sets}
\label{sec:niches}

Let $n$ index a \term{constructional niche} (a communicative job), and let
$\mathcal{V}_{n,t}^{\star}$ be the opportunity set: the attested,
hierarchically constructible, and analogically generated candidate realizations
that could plausibly do that job in some conditioning states. Membership in
this set doesn't presuppose that a candidate is licensed.
Let $N_t(n,c)$ be the number of opportunities for niche $n$ in conditioning
state $c$ during the current window $W_t$, and $k_t(x,n,c)$ the observed count of
candidate realization $x\in\mathcal{V}_{n,t}^{\star}$.

The star matters. Preempted gaps and innovations often involve candidates with no
positive tokens. Their counterfactual utilities are still defined because the
hierarchy $\mathcal{H}$ generates analogical candidates. English learners can
represent a left-branch extraction candidate by combining interrogative
fronting, modifier--noun packaging, and question formation, even if no one
licenses that candidate as an English resource. The same machinery lets a novel
but ordinary sentence inherit coverage from licensed higher-grain constructions
rather than being preempted merely because that exact sentence type is new.

\subsection{Usage as \texorpdfstring{licensing $\times$ choice}{licensing x choice} among candidates}
\label{sec:usage}

Separate \term{licensing} from \term{selection}. A candidate can be licensed but
rarely chosen.

Let $\rho_t^\star(x\mid n,c)\in[0,1]$ be the candidate-conditional
counterfactual probability of choosing $x$ if every registered candidate were
available in niche $n$ under $c$. A flexible choice model is a softmax over
utilities:
\[
\rho_t^\star(x\mid n,c)=
\frac{e^{U_t(x;n,c)}}{\sum_{x'\in\mathcal{V}_{n,t}^{\star}}e^{U_t(x';n,c)}},
\qquad
U_t(x;n,c)=\mathbf{w}^{\top}\mathbf{f}(x;n,c).
\]
Because the production model also contains an outside option, the quantity that
enters a licensing--usage factorization is instead the full-choice share
\[
\widetilde{\rho}_t^\star(x\mid n,c)=
\frac{e^{U_t(x;n,c)}}
{e^{U_t(\bot;n,c)}+
 \sum_{x'\in\mathcal{V}_{n,t}^{\star}}e^{U_t(x';n,c)}}.
\]
The two coincide only when the outside option has negligible mass. Empirical
norming has to offer the relevant avoidance, periphrasis, or
repair-in-progress response as well as the candidate realizations; a
candidate-only forced choice estimates $\rho_t^\star$, not
$\widetilde{\rho}_t^\star$.

At the population level, the expected usage rate is
\[
\pi_t(x\mid n,c)=
E_i\!\left[
\frac{z_i(x,c)e^{U_t(x;n,c)}}
{e^{U_t(\bot;n,c)}+
\sum_{x'\in\mathcal{V}_{n,t}^{\star}}z_i(x',c)e^{U_t(x';n,c)}}
\right],
\]
where, for a candidate $x$ with form--value pair $(f_x,v_x)$,
\[
z_i(x,c)=\mathbf{1}\!\left[
\exists A\in\mathcal{A}(f_x,v_x):
\operatorname{def}(A)\wedge\operatorname{sat}_t(A,c)\wedge
\bigwedge_{\kappa\in A}z_{i,t}(\kappa,c)
\right]
\]
abbreviates speaker $i$'s availability of a complete licensed route for $x$
(with the time index suppressed on the left), and $\bot$ is the outside option
(avoidance, periphrasis, or repair-in-progress). A candidate realization appears
once in the softmax: if several assemblies license the same realization, the
existential above aggregates them into the single gate $z_i(x,c)$ rather than
adding separate choice masses. With several required nodes,
the expectation is evaluated over the joint inclusion model rather than inferred
from marginal node means alone. When one pivotal availability gate dominates
and the other gates and full-set choice weights have been normed independently,
the predictive product
$\pi_t(x\mid n,c)\approx S_t^\theta(x,c)
\widetilde{\rho}_t^\star(x\mid n,c)$ is a useful population-level
approximation. In single-pivotal-node cases this reduces to
$\theta_t(\kappa_x,c)\widetilde{\rho}_t^\star(x\mid n,c)$, where $\kappa_x$ is
the node that would realize $x$ in $c$. The corresponding analyst prediction
replaces $S_t^\theta$ with $G_t$, or $\theta_t$ with $\widehat C_t$, as an
estimation step; it doesn't identify the ontic usage rate with the estimate.
The full-choice term can be high even when the population availability is near
zero, so absence becomes informative under that condition.

\subsection{Dynamics: evidence streams, weighted omissions, and repair}
\label{sec:update}

Speaker $i$ represents each construction--situation pair $(\kappa,c)$ with a
Beta filter over the population licensing rate $\theta_t(\kappa,c)$
(\S\ref{sec:objects}); the analyst may use the same family on a different data
stream. The Beta is exact after direct Bernoulli inclusion observations, but the
observable data here are gated choices. A target choice supplies a factor
proportional to $\theta$; a non-target choice supplies an affine factor in
$\theta$. The resulting finite-window posterior is generally a Beta mixture.
The bounded filter below projects that exact mixture back to Beta by matching
its first two moments.

Discount accumulated evidence before applying the new window. For baseline
prior $\mathrm{Beta}(a^0,b^0)$ and memory $\delta_m\in(0,1]$, put
\begin{equation}
\widetilde a_t=a^0+\delta_m(a_t-a^0),\qquad
\widetilde b_t=b^0+\delta_m(b_t-b^0).
\label{eq:baseline-discount}
\end{equation}
A quiet filter returns to its stated prior. The recurrence discounts the
accumulated state, and the current window enters once. Setting $\delta_m=1$
gives a fixed-parameter filter, while $\delta_m<1$ is an assumed-density dynamic
filter for a drifting convention \citep{WestHarrison1997}.

Consider a non-target outcome $y_j\ne x$. Let
\begin{equation}
\ell_j(x)=\log P(y_j\mid Z_x=0,n,c)-
\log P(y_j\mid Z_x=1,n,c)
\label{eq:attributability}
\end{equation}
be its omission log likelihood ratio, used here as an information diagnostic.
The update itself uses the probability-space contrast below. With $r_j\in[0,1]$
the independently coded probability that the occasion is attributable to the
target niche, define
\begin{equation}
d_j=r_j\left(1-
\frac{P(y_j\mid Z_x=1,n,c)}{P(y_j\mid Z_x=0,n,c)}\right)\in[0,1].
\label{eq:preemption-mass}
\end{equation}
Marginalizing the latent availability state
$Z_x\sim\mathrm{Bernoulli}(\theta)$ gives, up to a constant independent of
$\theta$,
\[
P(y_j\mid\theta,n,c)\ \propto\ 1-d_j\theta.
\]
Attribution scales $d_j$ in probability space: with probability $1-r_j$ the
event is uninformative, and with probability $r_j$ it carries the full
counterfactual contrast.

If the window contains $s_t$ target choices and $m_t$ non-target choices, its
exact posterior kernel after~\eqref{eq:baseline-discount} is
\begin{equation}
p_{t+1}(\theta)\ \propto\
\theta^{\widetilde a_t-1+s_t}(1-\theta)^{\widetilde b_t-1}
\prod_{j=1}^{m_t}(1-d_j\theta).
\label{eq:affine-window-posterior}
\end{equation}
The affine product has the Bernstein expansion
\[
\prod_{j=1}^{m_t}(1-d_j\theta)=
\sum_{k=0}^{m_t}c_k\theta^k(1-\theta)^{m_t-k},
\]
where $c^{(0)}_0=1$ and
$c^{(j+1)}_k=c^{(j)}_k+(1-d_{j+1})c^{(j)}_{k-1}$, with out-of-range
coefficients zero. Hence~\eqref{eq:affine-window-posterior} is exactly a mixture
of at most $m_t+1$ Betas,
\[
\mathrm{Beta}(\widetilde a_t+s_t+k,
              \widetilde b_t+m_t-k),
\]
with normalized weights proportional to the corresponding $c_k$ times the
Beta normalizer. Let $\mu_{t+1}$ and $v_{t+1}$ be that mixture's exact mean and
variance. The bounded assumed-density step is
\begin{align}
\nu_{t+1}&=\frac{\mu_{t+1}(1-\mu_{t+1})}{v_{t+1}}-1, &
a_{t+1}&=\mu_{t+1}\nu_{t+1}, &
b_{t+1}&=(1-\mu_{t+1})\nu_{t+1},\\
C_{t+1}&=\frac{a_{t+1}}{a_{t+1}+b_{t+1}}.
\label{eq:beta-update}
\end{align}
This projection preserves the exact first two posterior moments but not, in
general, higher moments. Its approximation error is part of the model
audit. Direct target observations, $d_j=0$, and $d_j=1$ recover the relevant
conjugate limits.

The equations describe one pivotal availability gate. If a token's parse or
niche attribution is uncertain, that uncertainty is marginalized in its
likelihood before $d_j$ is calculated. A multi-node target requires
the joint inclusion model from \S\ref{sec:objects}; one omission LLR can't be
split into independent node failure counts.

Population change enters through an explicit learner-to-population transition
kernel. Let $C^{(i)}_{\kappa,c,t}$ and $\nu^{(i)}_{\kappa,c,t}$ be speaker $i$'s
filter summaries, let $\xi_{i,t}$ collect lect, network, and other independently
specified adoption covariates, and let $\lambda_{i,\kappa}\in[0,1]$ govern
persistence. Then
\begin{equation}
z_{i,t+1}(\kappa,c)\sim\operatorname{Bernoulli}\!\left(
  (1-\lambda_{i,\kappa})z_{i,t}(\kappa,c)+
  \lambda_{i,\kappa}h_\kappa(C^{(i)}_{\kappa,c,t},
    \nu^{(i)}_{\kappa,c,t},\xi_{i,t},c)\right),
\label{eq:inclusion-transition}
\end{equation}
where $h_\kappa\in[0,1]$ is an adoption--retention response fixed independently
of the outcome being predicted. The numerical analysis below uses
\begin{equation}
h_\kappa(C,\nu;q_\kappa)=
P(\theta\geq q_\kappa\mid C,\nu),
\label{eq:posterior-threshold-adoption}
\end{equation}
which can be read as posterior sampling against an adoption criterion.
$q_\kappa$ has to be fixed from independently measured switching or coordination
costs, while concentration supplies the response's curvature. This causal
transition parameter is distinct from the judge's non-constitutive read-out
threshold $\tau(c)$ in \S\ref{sec:decision-readout}. A logistic
response to $C-q_\kappa$ remains a reduced-form sensitivity analysis only when
its slope is independently estimated. The persistence term allows stable individual
lects even when the population is divided. For a finite population of $M$
speakers, distinguish the empirical prevalence
\[
\theta^{(M)}_{t+1}(\kappa,c)=\frac{1}{M}\sum_{i=1}^{M}
z_{i,t+1}(\kappa,c)
\]
from the exchangeable rate $\theta_{t+1}$. Conditional on the current population
state, the expected empirical prevalence is the average of the transition
probabilities in~\eqref{eq:inclusion-transition}; in a homogeneous mean-field
case the realized prevalence fluctuates around that common probability at
$O_p(M^{-1/2})$. Replacing the average of a nonlinear $h_\kappa$ by its value at
average covariates is a separate closure approximation whose error depends on
population heterogeneity, not on the opportunity count $N_t$. Speakers'
evidence states can mediate population change, while an analyst's
posterior remains only a measurement of that change.
\noindent\textsc{Counterfactual informativeness.} When $y_i$ is an
in-repertoire competitor, let
\[
D_{0,i}=e^{U_t(\bot;n,c)}+
  \sum_{x'\ne x}z_i(x',c)e^{U_t(x';n,c)}
\]
be the gated normalizer when $x$ is unavailable, and define the occasion-level
full-choice share
\[
\widetilde{\rho}_{i,t}^\star(x\mid n,c)=
  \frac{e^{U_t(x;n,c)}}{D_{0,i}+e^{U_t(x;n,c)}}.
\]
Holding the other gates fixed, normalized softmax choice then gives the exact
identity
\[
\frac{P(y_i\mid Z_x=1,n,c)}{P(y_i\mid Z_x=0,n,c)}
  =1-\widetilde{\rho}_{i,t}^\star(x\mid n,c),
\]
and hence
\[
\ell^{\mathrm{omit}}_{i,t}(x)
=-\log\!\bigl(1-\widetilde{\rho}_{i,t}^\star(x\mid n,c)\bigr)
  \approx \widetilde{\rho}_{i,t}^\star(x\mid n,c)
\]
only when the target's full gated mass is small. The common
$N_t\cdot\rho_t^\star$ score is recovered as a small-mass information
approximation in a near-complete gap and
under four further conditions: nearly every opportunity is an attributable
omission ($r_i\approx1$), the outside option has negligible mass, the other
candidate gates match the norming task, and the target mass is small enough for
the log expansion. Continued attributable competitor choices can keep
supplying evidence against a concentrated gap; reopening it requires an
actuation route that changes utilities, gates, or the observed likelihood.

When $y_i$ is the outside option--periphrasis or avoidance, available at fixed
utility under both hypotheses--the same denominator identity applies. Its
choice is nearly uninformative only when the target's full counterfactual mass
$\widetilde{\rho}_{i,t}^\star$ is small, which requires $x$'s utility to be low
relative to the candidate set and the outside option. That condition is
independently estimable from the same norming task. Then
$\ell^{\mathrm{omit}}_{i,t}(x)\approx0$.

When candidate support is weak, the affine contrast lets periphrasis remain
weak evidence against any one candidate. This separates the defective cell's
low-concentration state from the preempted gap's high-concentration state
(\S\ref{sec:winnerless-cells},
\S\ref{sec:worked-examples}).

\noindent\textsc{The repair stream: a candidate controller.}
\label{sec:repair-controller}
The paper's support claims come in grades: stabilization, maintenance, control
(\S\ref{M-sec:implications}). The transition and full map below specify
candidate conditions for stabilization; showing that it or the proposed loop
supports an actual linguistic profile requires longitudinal or intervention
evidence. Control needs the stronger result that a process registers departures
from the licensed profile and pushes the population back. Conversational repair
is the candidate considered here.

Let $\Delta(d)$ be the expected common-ground divergence from mis-setting
operator dimension $d$: the quantitative face of \enquote{update-configuring}
in the working test of \S\ref{M-sec:operator-values}. Formally, $\Delta(d)$ is an
expected KL divergence between the common-ground posterior induced by the
intended value and the posterior induced by the mis-set value, in the same
spirit as RSA-style update models \citep{GoodmanFrank2016}. When a token
mis-sets $d$, update-oriented repair occurs with probability
$r_{\mathrm{upd}}(\Delta,\iota)$, increasing in $\Delta$ and modulated by
footing $\iota$. A different stream targets a value-intact exponent or
construction whose licensing evidence has concentration $\nu$; write its
probability as $r_{\mathrm{form}}(\nu,\iota)$. Over a window with
$N_t(n,c)$ opportunities, the two candidate repair exposures are
\begin{equation}
\begin{aligned}
\psi_{\mathrm{upd}}(n)
  &=N_t(n,c)\,P(\text{wrong value})\,
    r_{\mathrm{upd}}(\Delta,\iota),\\
\psi_{\mathrm{form}}(n)
  &=N_t(n,c)\,P(\text{value intact, form excluded})\,
    r_{\mathrm{form}}(\nu,\iota).
\end{aligned}
\label{eq:repair-flow}
\end{equation}
The first magnitude increases with divergence through $r_{\mathrm{upd}}$; the
second can remain high when $\Delta\approx0$ because the intended update is
intact while the form is strongly excluded. These are candidate repair
exposures, not yet derived terms in the likelihood-consistent filter.

The numerical results in \S\ref{sec:emergent-categoricality} supply no evidence
for the repair hypothesis; the companion software's programmed repair
comparison is only a diagnostic. Repair earns controller status under
Woodward's interventionist criterion if an
intervention on it changes the profile's return-to-attractor behaviour in an
invariant, predictable way \citep{woodward2001,
woodward2003MakingThingsHappen}. Until then, repair remains a testable candidate
controller.

Two commitments follow. First, the account predicts a distribution over repair
targets, not a binary grammar/style split. Opportunity is estimable from
opportunity-annotated corpora; $\Delta$ comes from comprehension probes fixed
before the repair data is inspected; and $\nu$ comes from the independent
licensing evidence. A closed, dense contrast with $\Delta\approx0$ may still be
categorically excluded through constructional or exponent licensing, and may
attract form correction. What it should lack is the update-oriented
repair profile predicted by $r_{\mathrm{upd}}$. Second, if that repair profile
tracks social indexing rather than $\Delta$ after the other factors are held
fixed, the proposed operator-directed controller fails cleanly. Blythe and
Croft's utterance-selection dynamics remain the no-repair special case, so the
import is conservative \autocite{BlytheCroft2012}.

\noindent\textsc{Grain asymmetry and compositional inheritance.} Positive and
negative evidence can bear on different constructional grains. A heard token
supports the activated constructions that sanction its observed properties.
An attributable omission instead bears on the complete candidate-level route
that would have made the missing realization available. In Bayesian terms, the
likelihood penalty for structured non-occurrence falls on a hypothesis that
wastes choice probability on unobserved sentence types
\autocite{PerforsEtAl2011}.

Compositional inheritance raises an objection. Token-level licensing is assembled from
constructional nodes, so left-branch extraction activates the well-licensed
nodes for interrogative fronting, question formation, and modifier--noun
packaging. Their independent support might then seem to make the discontinuous
candidate available.

That result doesn't follow from the status definition. Availability requires a
complete licensed route, and those parent nodes sanction only \emph{contiguous}
realizations. None of them covers the
determiner--head discontinuity that defines
\ungram{\mention{Which did you buy car?}}; producing that token requires a node
that licenses a determiner--head arc crossing the verb, and no
contiguity-preserving node in $\mathcal{H}$ generates one
\autocite{reynolds2026lbe}.

The parents can retain high licensing on the strength of their abundant
contiguous tokens without licensing the discontinuous route. If the grammar
posits a dedicated left-branch-extraction node, its licensing has to be
estimated separately. Attributable omissions then bear on the candidate-level
gate only to the extent that competitors would have been chosen instead, which
is why the independent $\widetilde{\rho}_t^\star$ norming required in
\S\ref{sec:identification-c} is necessary.
This asymmetry also answers an acquisition worry. If token-level licensing were
just pooled parent licensing, learners should pass through a stage of treating
left-branch extraction as live. They don't need to: the candidate is generated
by analogy, so it's on the table with non-zero prior support, but its dedicated
node has no scaffolding of its own, fronting opportunities are dense in the
input, and each one adds preemption mass.

The quantitative model in this supplement updates a candidate-level pivotal
gate. It doesn't yet define fractional evidence allocation across all activated
nodes or a numerical pooling rule. Evidence distribution through a construction
hierarchy remains a further modelling problem; the equations here don't settle it.

A companion corpus study quantifies
one side of the evidential situation: across 1{,}228 fronting opportunities in a
sixteen-corpus Universal Dependencies sample, zero genuine determiner--head
discontinuities occur (95\% upper bound 0.24\%), while contiguity-preserving
alternatives (pied-piping, fused-head construals, and the \mention{big mess}
family) occur at measurable rates \autocite{reynolds2026lbe}. Comprehension
pushes the same way: a fronted \mention{whose} or \mention{which} defaults to a
fused-head parse, so the discontinuous analysis also loses the parsing race
\autocite{Culicover_Varaschin_Winkler_2022_Radical}. What remains to be normed
is $\widetilde{\rho}_t^\star$: if speakers would rarely choose the discontinuous form even
under a permissive stipulation that English allowed it, the absence would support
starvation rather than preemption. The predicted non-satiation of left-branch
extraction remains a test; \posscite{Lu2024} meta-analysis of
satiation in extraction from islands supplies the closest current analogue, not
a direct measurement of \abbr{lbe} itself.

\noindent\textsc{Open-loop learning and closed-loop dynamics.} The discrete
update~\eqref{eq:beta-update} changes an evidence state in a fixed exogenous
environment. Figure~\ref{fig:posterior-means} shows the resulting evidential
difference between sparse and dense attributable omissions. Population
bifurcation requires the transition in~\eqref{eq:inclusion-transition}:
inclusion gates production, production supplies later choice evidence, and
speaker filters affect later adoption and retention.

\subsection{Conditional emergence of separated regimes}
\label{sec:emergent-categoricality}

The sense of \term{emergence} used here is weak and mechanistic:
population-level categoricality is a possible outcome of learner updating,
adoption, and selection. This differs from Hopper's \enquote{emergent grammar}, where grammar is
perpetually provisional \autocite{hopper1987}.

For a homogeneous focal candidate, let
$X_t=(\theta_t,a_t,b_t)$ and let
$\mathcal M_{\theta_t}(a_t,b_t)=(a'_t,b'_t)$ be the deterministic
expected-window version of the discounted likelihood update in
\eqref{eq:affine-window-posterior}. It uses the expected target and non-target
counts implied by the gated-choice model and moment-matches the resulting
power-likelihood posterior. The composed mean-field map is
\begin{equation}
F(X_t)=\left(
(1-\lambda)\theta_t+\lambda h_\kappa(a'_t,b'_t),
a'_t,b'_t\right).
\label{eq:full-mean-field-map}
\end{equation}
This expected-window closure is a deterministic diagnostic of the
finite-population stochastic process, whose transition kernel remains
\eqref{eq:inclusion-transition}. Fixed points satisfy
$F(X^\ast)=X^\ast$; their local stability is governed by the full Jacobian,
\begin{equation}
\rho\!\left(J_F(X^\ast)\right)<1,
\label{eq:full-map-stability}
\end{equation}
where $\rho$ is spectral radius. This criterion retains the slow evidence
directions that a scalar population equation discards.

A reduced crossing plot remains useful for location. Let $\eta$ be expected
opportunity inflow per step before observation and let $q$ be the adoption
threshold. Holding $\theta$ fixed, let
$(a^\ast_{\eta}(\theta),b^\ast_{\eta}(\theta))$ be the stationary evidence state
and put
\[
H_{\eta,q}(\theta)=
h_{\kappa,q}(a^\ast_{\eta}(\theta),b^\ast_{\eta}(\theta)),
\qquad
R(\theta;\eta,q)=H_{\eta,q}(\theta)-\theta.
\]
Here $h_{\kappa,q}(a,b)$ abbreviates
$h_\kappa(a/(a+b),a+b;q)$ from
\eqref{eq:posterior-threshold-adoption}.
Zeros of $R$ locate candidate equilibria. They do
not supply return times: the scalar derivative can differ sharply from the
dominant eigenvalue of~\eqref{eq:full-map-stability} when evidence is slow.
The unstable interior crossing is denoted $\theta^\ddagger$ below.

The numerical analysis uses a specified two-candidate cell: equal candidate
utilities, outside utility $-4$, twelve opportunities per step, observation
probability $.6$, memory $.96$, inclusion-update rate $.08$, a uniform Beta
baseline, and $q_\kappa=.5$. Under the posterior-threshold response in
\eqref{eq:posterior-threshold-adoption}, the deterministic closure has stable
fixed points effectively at $0$ and $1$ and an unstable crossing at
$\theta^\ddagger\approx .493$. The corresponding spectral radii are $.961$,
$1.087$, and $.971$. A logistic reduced form with slope $8$ has stable fixed
points near $.024$ and $.968$ around an unstable crossing near $.489$; slope
$4$ has only one stable interior fixed point. A proportional response likewise
has one interior fixed point. Thus the independently specified adoption response
supplies the nonlinearity.

Finite-population runs preserve that scope result. Across 36 seeded runs of 320
steps initialized at $.1$, $.5$, and $.9$, the posterior-threshold and
slope-$8$ responses both passed the declared basin-separation and endpoint
criteria; the proportional response didn't. A dominant outside option
($U(\bot)=5$) produced weak evidence and no endpoint concentration. Moderate
inclusion-only and evidence-state perturbations returned to both attracting
regions. These runs establish existence and scope for the declared process.
They provide neither an empirical fit nor a stationary-distribution result. The
exact grid, seeds, approximation diagnostics, and artifact hash are archived
with the companion software.

\noindent\textsc{Local bifurcation diagnostic.} The full expected-window closure
has a codimension-two cusp of fixed points if
\begin{equation}
R=\partial_\theta R=\partial^2_\theta R=0,
\qquad
\partial^3_\theta R\neq0,
\qquad
\det
\begin{pmatrix}
\partial_\eta R & \partial_q R \\
\partial_{\theta\eta}R & \partial_{\theta q}R
\end{pmatrix}
\neq0,
\label{eq:cusp-conditions}
\end{equation}
provided the full map has one unit eigenvalue and stable transverse
eigenvalues. These are local fixed-point conditions; they neither introduce a
potential nor invoke a global catastrophe classification.

Numerical continuation under the utilities, observation probability, memory,
and inclusion-update rate specified above, while varying $\eta$ and $q$, locates
the point
\[
(\theta_c,\eta_c,q_c)=(.370957,.479156,.552048),
\qquad
(a_c,b_c)=(1.798418,2.039839).
\]
The effective observed inflow is $.287494$. At the reported differentiation
scale, the first two state derivatives of $R$ are below $3\times10^{-8}$ in
absolute value, $\partial^3_\theta R=-7.2843$, and the unfolding determinant
is $2.3345$. The full-map eigenvalues are $1.000000031$, $.961496$, and
$.883699$; step-size audits preserve the non-zero third derivative and
determinant.

A nearby point inside the cusp wedge has fixed points at
$\theta\approx.253,.370,.505$, with stable outer points and an unstable middle
point; a point outside has one stable fixed point near $.230$. Five stationary
evidence solves from distinct initial states agree to within $10^{-8}$.

A quasi-static two-control cycle also distinguishes this geometry from a steep
single-valued response. Holding one local control coordinate fixed and sweeping
the other produces two static folds, at local biases $-.001041$ and $.000901$.
As the dwell time at each control value rises from 1{,}000 to 10{,}000 map
iterations, the directional jumps converge toward those folds.

The same current external controls can support different stable licensing
states when the retained evidence state differs. An empirical application would
have to estimate effective exposure and the adoption threshold without fitting
them to the transition points being explained.

\noindent\textsc{Slowing and its failed stochastic proxy.} At $q=.5$, the
lower-branch fold occurs at $\eta=.654741$. As the deterministic state approaches
that fold from $2.5$ to $1.02$ times the fold inflow, the dominant spectral
radius rises from $.964802$ to $.995641$, its implied relaxation time rises from
28 to 229 windows, and measured pulse recovery rises from 82 to 288 windows.
This establishes deterministic slowing in the expected closure.

It doesn't establish a usable stochastic early-warning signal. In the
finite-population diagnostic, each of four fold distances had 24 runs with 240
speakers over 3{,}000 steps. Median detrended variance and lag-one
autocorrelation rose through the $1.1$ condition, but 23 of 24 paths escaped
there. At $1.02$, all paths escaped, only four supplied a usable pre-escape
segment, and both summaries fell. The preregistered monotonicity and
minimum-replicate gates fail.

That diagnostic also uses one opportunity per step with observation probability
set to effective inflow; other opportunity--observation factorizations remain
untested. The local
bifurcation and deterministic hysteresis survive this failure, but no general
observable early-warning prediction follows.

The convergence claim has to match both levels of the model. With perfect
retention ($\delta_m=1$), an individual evidence filter is close to
reinforcement processes of Pólya type, where stochastic approximation supplies
relevant techniques \citep{pemantle2007}. The signaling-game case,
structurally the closest analogue, is treated directly by Argiento, Pemantle,
Skyrms, and Volkov \citep{argiento2009}, with an extension by Hu, Skyrms,
and Tarrès \citep{huSkyrmsTarres2011}. Those results don't establish
convergence of the composed population transition
in~\eqref{eq:inclusion-transition}; that requires a specified $h_\kappa$.

With bounded memory ($\delta_m<1$), the learner filters have constant gain, so
the decreasing-gain analogues no longer prove almost-sure convergence even at
that level. The population-level formal target is a stationary distribution for
the composed transition concentrated near the full map's attracting regions,
with escape times and residual dispersion set by opportunity, persistence, and
memory: high-opportunity contrasts should sit tightly near the extremes, while
low-opportunity or outside-option-dominated contrasts remain diffuse. That claim is a proof
obligation rather than a theorem here.

The operator-specific prediction also needs its scope stated. Neither the full
map nor the local bifurcation diagnostic derives the operator-specific response
boundary. The OVMG claim is
conditional: adoption criteria and selection intensity have to vary with
independently estimated switching or coordination costs, opportunity, and
update divergence. Under that condition, dense, update-configuring contrasts
may concentrate faster.

Low $\Delta$ doesn't by itself imply gradient
licensing: a construction or exponent can be categorically excluded by dense
competitor evidence without being an operator. The operator-specific claim is
instead about comprehension and repair type. If a closed, dense,
non-configuring contrast shows the same update-confusion and update-oriented
repair profile as a matched operator contrast, that part of the account fails.

Bounded memory adds a quantitative residual: even near an attractor, dispersion
should scale with opportunity and forgetting, roughly with the effective
concentration supplied by the discounted evidence filter.

Sampling drift near the extremes makes the boundaries sticky in finite
populations, while middle states stay vulnerable to renewed evidence. The same
escape-time logic connects winnerless-cell metastability (\S\ref{sec:winnerless-cells})
and actuation (\S\ref{sec:actuation}): a state can persist for long periods and
then move quickly once the observation likelihood or outside-option utility changes.
This connects the dynamics to complex-adaptive and cultural-evolutionary
accounts of language change \autocite{beckner2009,kirby2014,BlytheCroft2012},
and to invisible-hand explanations, where macro-level order emerges as the
unintended consequence of micro-level choices \autocite{keller1994}.

\subsection{Derived obligatoriness}
\label{sec:derived-obligatoriness}

The state theory defines saturation relative to the obligatory dimensions
$\mathrm{OBL}_t(n,c,A)$ (\S\ref{sec:objects}), where $n=n(A,c)$ is the
registered at-issue niche. Those dimensions are path-dependent population
properties, estimated separately by an analyst. Let
$\mathcal{A}_0(d,n,\phi,c)$ be the finite, independently fixed class of defined
assemblies of frame type $\phi$ that realize niche $n$ while leaving $d$ unset.
For $A$ in that class, define the population's \term{raw zero-marking rate}
\[
R_s^0(A,n,c)=P_j\!\left(\text{speaker $j$ realizes a $d$-at-issue opportunity
in $(n,c)$ with $A$ at $s$}\right).
\]
This realized population statistic is distinct from the licensing predicate
$L_s$ and from a status estimate. The same simple-present nodes can have high
licensing in a habitual niche while their zero-marked realization rate is near
zero in an ongoing-at-issue niche. General node availability, niche-specific
zero marking, and saturation remain separate quantities.

Let $W_t=[t-W,t)$ be a preceding convention-individuating window, and let
$\Pi_{W_t}(A,d,n,c)$ be the realized pressure against $A$ there: $d$-at-issue
opportunities in which a marked competitor wins, weighted by the independently
normed full-choice share $\widetilde\rho^\star$. This weight is ungated; it
doesn't use $\operatorname{sat}_s$ or $\mathrm{OBL}_s$. Outside-option choices
don't contribute. Write $E_t(d,n,\phi,c)$ for the entry condition that, for
every $A\in\mathcal A_0(d,n,\phi,c)$,
\[
\sup_{s\in W_t}R_s^0(A,n,c)<\epsilon_0
\qquad\text{and}\qquad
\Pi_{W_t}(A,d,n,c)\geq\pi_0.
\]
Write $X_t(d,n,\phi,c)$ for the release condition that, for some such $A$,
\[
\inf_{s\in W_t}R_s^0(A,n,c)>\epsilon_1
\qquad\text{or}\qquad
\Pi_{W_t}(A,d,n,c)\leq\pi_1,
\]
where $\epsilon_1>\epsilon_0$ and $\pi_1<\pi_0$. Given an initial convention
state $\mathrm{OBL}_{t_0}$ fixed independently of the target interval, the
path-dependent macro is
\[
d\in\mathrm{OBL}_t(n,c,\phi)
\quad\Longleftrightarrow\quad
E_t(d,n,\phi,c)\ \lor\
\bigl[d\in\mathrm{OBL}_{t-1}(n,c,\phi)\land
\neg X_t(d,n,\phi,c)\bigr].
\]

The entry clauses say that every zero-marking realization has persistently
disappeared under dense, attributable competition rather than because the niche
was starved or an outside option absorbed a winnerless cell. The release clauses
require sustained recovery or a clear pressure lapse. The two pairs of
thresholds leave a hysteresis band in which the preceding state persists. The
timescale and thresholds are
community-indexed model parameters fixed independently of the contrast being
classified. The initial condition and strict temporal descent make the
recursion well founded: realized outputs and ungated pressure in $W_t$, together
with $\mathrm{OBL}_{t-1}$, determine $\mathrm{OBL}_t$;
$\mathrm{OBL}_t$ determines $\operatorname{sat}_t$; and saturation then enters
$S_t^\theta$. This makes status a functional of the recent population
history, not just of instantaneous $\boldsymbol{\theta}_t$.

An analyst identifies $d$ as obligatory when the posterior upper bound for every relevant
$R_s^0(A,n,c)$ lies below $\epsilon_0$, the estimated pressure clears $\pi_0$, and the
effective evidence meets a preregistered $\nu_{\min}$. Those conditions warrant
an inference about the population property $\mathrm{OBL}_t$. Identification of
release analogously requires posterior evidence for $X_t$ across the registered
window.

The proposed mechanism is the population transition plus preemption applied to
zero-marking. In a $d$-at-issue niche, marked and unmarked assemblies compete;
where the marked assembly repeatedly wins, speakers' filters and later inclusion
states can drive the raw zero-marking rate down. The full map shows one
possible phase portrait for that process, while the stationary-distribution
claim remains a proof obligation rather than part of the definition of
obligatoriness.

The typological payoff is that \S\ref{M-sec:community-values}'s descriptions
convert to predictions. On the proposed diagnosis, English progressive marking
is obligatory because simple-present assemblies in ongoing-at-issue niches have
persistently low raw zero-marking rates under a winning marked competitor; French
zero-marking there hasn't, so the dimension stays optional, and
\mention{Elle étudie maintenant} is fine. Turkish evidential marking is the
same kind of hypothesis on a different dimension. Grammaticalization's
obligatorification receives a candidate mechanism: a contrast on its way
to obligatory status is one whose zero-marking competitor is losing dense
niches one at a time, which predicts an intermediate stage--obligatory in some
frame types and niches, optional in others--that grammaticalization corpora are
positioned to test.

\subsection{Winnerless cells}
\label{sec:winnerless-cells}

Preempted gaps have a winning competitor. Defective paradigm cells don't, and
the dynamics explain how the difference can be stable rather than transitional.

Suppose candidates $f_1,\dots,f_m$ for a dense cell have weak analogical support,
the outside option (periphrasis, avoidance) has moderate fixed utility, and
producing a weakly licensed form carries a face cost. Then the production rule
routes the cell's opportunities to the outside option: no candidate accumulates
positive evidence. The omission likelihood matters here: outside-option choices
have $\ell^{\mathrm{omit}}_{i,t}(x)\approx0$, so no candidate accumulates much negative evidence
either. Evidence starvation alone preserves the prior. The
low-mean/low-concentration profile requires an
additional independently motivated ingredient--low priors for unsupported
analogical candidates, weak diffuse negative evidence, a hierarchical prior over
defective cells, or a latent cell-level state such as \enquote{no candidate is
established}. Given such an ingredient, avoidance keeps every candidate at low
mean and low concentration, because escaping the state requires some candidate
to accumulate positive evidence that the outside option keeps starving it of.

The claim has two strengths. The winnerless-cell example in
\S\ref{sec:worked-examples} relies on the qualitative mechanism: an avoidance
attractor plus evidence starvation, with low prior support or weak diffuse
negative evidence supplying the low mean. The stronger metastability claim,
persistence for times exponential in population size, remains a conjecture with
analogues in evolutionary dynamics.
\posscite{Sims2023Defectiveness} point that the maintenance mixture varies
across cases--structural uncertainty, lexical restriction, learning dynamics,
avoidance--is compatible with this narrower claim: the model supplies the
avoidance-and-starvation component and is agnostic about the rest of the
mixture.

The phenomenological contrast is represented by the resulting evidence profile.
Clean defective cells
(\mention{pobedit'} \abbr{1sg}, with \ungram{\mention{amn't}} depending on niche
granularity) show low mean with
low concentration--\enquote{I don't know how to say it}--where preempted gaps
show low mean with high concentration--\enquote{that's impossible}. Same mean,
different concentration, different feeling, read out by $\Phi$
(\S\ref{sec:feeling-new}). The satiation-framing prediction follows: defective
cells should behave like low-evidence forms, movable under framing; preempted
gaps shouldn't \autocite{Broadbent2009Amnt,Pertsova2016RussianGaps,
ThomsEtAl2025Amnt}.

\subsection{Destabilization of moribund contrasts}
\label{sec:destabilization}

Bounded memory yields one more result. With discount $\delta_m<1$, the
equilibrium concentration of a node's posterior is set by the balance of
evidence inflow against forgetting:
$\nu^{\ast}\approx\nu_0+O\bigl(N_t(n,c)/(1-\delta_m)\bigr)$, with $N_t(n,c)$ the
per-window opportunity count for the node's niche and $\nu_0=a^0+b^0$ the baseline
concentration. If opportunity falls--a case distinction going moribund, a niche
drying up--equilibrium concentration falls toward the baseline. With no further
evidence, means drift toward the baseline prior under~\eqref{eq:beta-update};
with sparse, idiosyncratic evidence, individual histories can diverge. Falling
excess concentration is the direct consequence and the primary preregistrable
precursor; increasing test--retest instability is a behavioral proxy only if it
is independently validated against concentration.

Prior heterogeneity alone doesn't make dispersion lead the mean. If, after a
shared evidence stream disappears,
\[
C_{i,t}=w_t m_i+(1-w_t)C_E,
\]
where $m_i$ is speaker $i$'s baseline-prior mean, $C_E$ is the former common
evidence mean, and $w_t$ rises toward one, movement of the population mean is
linear in $w_t$ while between-speaker variance is
$w_t^2\operatorname{Var}(m_i)$. On a common fraction-of-eventual-change scale,
dispersion lags. The companion simulation reproduces this ordering at
both 25\% and 50\% progress. Dispersion-leading decline remains a conditional
hypothesis requiring a faster heterogeneity source, such as cohort replacement,
unequal network opportunity loss, heterogeneous retention, or a specified
judgment mapping.

\subsection{Actuation}
\label{sec:actuation}

In a bistable cell of the full map, \term{actuation} is movement of the coupled
population-and-evidence state across the basin boundary whose population
coordinate is $\theta^\ddagger$. A utility change, new target evidence, altered
outside-option mass, or a social change in $q_\kappa$ can move the state across
that boundary. Speakers' filters then alter adoption and retention through
\eqref{eq:inclusion-transition}. An analyst posterior should track the
population movement with a lag; motion in its mean isn't itself actuation.
No S-shaped time course follows without specifying the full perturbation and
the slow evidence state.

The factors that drive such bifurcations align with the motivations discussed
in \S\ref{M-sec:motivations}. Semantic reanalysis may increase utility (and thus
$\widetilde{\rho}_t^\star$) by making a form--value relation easier to interpret or more useful.
Social pressures may shift utilities (prestige effects) or alter the relevant
community standard. Structural motivations such as analogical extension can
systematically raise licensing probability by borrowing strength from frequent
neighbors. Processing innovations may reduce error rates or locality
costs, changing either the observation likelihood or later adoption.

Individual innovation is insufficient. A few speakers adopting a marginal
construction don't guarantee its success; the likelihood and transition have to
shift so that inclusion is systematically favoured across the relevant
community. Many sensible innovations fail because they appear before those
community conditions are in place.

Near the separatrix, small perturbations in community attitudes can trigger
rapid shifts between exclusion and licensing. As the state moves further past
$\theta^\ddagger$, reversal becomes increasingly unlikely. The full map can
represent the S-curves documented in some language changes
\autocite{kroch1989}. For constructions in the marginal zone, with the
population state near $\theta^\ddagger$, the model predicts
heightened sensitivity to external factors and greater cross-community
variation. Small or peripheral communities may show volatile behaviour near
bifurcation points before a change spreads to larger population centres.

The destabilization result sharpens the actuation picture at the other end of a
form's life. Actuation is crossing a separatrix; destabilization is a collapse
in concentration. A contrast can lose its grip when evidence of any kind stops
arriving, without accumulating negative evidence; the model predicts the two
routes are empirically
distinguishable: actuation crosses a basin boundary in the coupled state;
moribundity first removes evidence confidence and need not cross that boundary
until later population turnover. Judgment dispersion may lead only under the
additional heterogeneity mechanisms named above.

\subsection{Apparent exclusions from concentrated non-licensing}
\label{sec:categorical-without-bans}

When a defined, saturated assembly exists but its pivotal node's licensing
posterior is driven toward zero and concentrates there under persistent
preemption in a niche with many observed opportunities, speakers converge on a sharp,
non-satiating rejection profile. Additional processing penalties in $R_i(u,c)$
(\eg systematic garden-path repair due to an entrenched fused-head construal)
can strengthen the subjective categoricality without any independent structural
constraint being posited.
